% =====================================================================
% CHAPTER 1: Introduction  (~7 pages)
% WRITING SESSION INSTRUCTIONS (keep until the chapter is done):
%   - Read thesis/WRITING_PLAN.md section 3 FIRST (content spec, sources,
%     figures, PO evidence per section).
%   - Invoke ruet-thesis-expert (format) + research-linguistics-expert
%     (prose). Register: 4-paragraph funnel stretched to chapter length
%     (patterns.md section 2); assert, do not hedge.
%   - Every number cites a repo source; no en/em dashes in prose; log PO
%     evidence lines to thesis/book/po_tracker_notes.md as you draft.
%   - POs demonstrated here: PO1 (engineering knowledge), PO2 (problem
%     analysis). CEP K-profile (K3, K4, K5, K8) stated in section 1.4.
% =====================================================================
\chapter{Introduction}

\section{Background and Motivation}
% Sources: review_paper/sections/01_introduction.tex + 02_background.tex,
% papers/summaries/tier1/surveys/Tosi_Survey_IJCV2025.md.
% Content: why machine depth perception matters; passive stereo vs LiDAR
% vs structured light (cost, power, form factor); state Eq 1.1 EARLY with
% a where-clause; applications (mobile robots, drones, AR, embedded rigs).
Depth perception is a prerequisite for a machine that must act in
three-dimensional space. A mobile robot avoids obstacles, a drone holds
its position, and an augmented-reality headset anchors virtual content
only if it knows how far away each surface lies. The sensing options differ sharply in cost and character. Scanning LiDAR
measures time of flight along emitted laser beams and delivers accurate
but sparse points at a price and power draw that exclude small platforms.
Structured-light and time-of-flight cameras emit infrared patterns or
pulses and recover dense depth indoors, but their emitters limit range,
fail in sunlight, and interfere with one another when several devices
share a space. Among the passive ways to recover depth, binocular stereo
is attractive for embedded systems because it emits nothing, draws little
power, works indoors and outdoors alike, and uses two inexpensive image
sensors rather than a laser scanner; its cost is computational, since the
correspondence between the two views must be solved for every pixel, and
that computation is precisely what this thesis makes affordable. Given a pair of images from two horizontally
separated cameras, the horizontal shift of a scene point between the two
views, its disparity, determines depth through the calibrated relation of
\eqref{eq:depth_from_disparity}. A dense disparity map therefore yields a
dense depth map, which downstream perception and reconstruction consume
directly. The applications that motivate this thesis share a common
trait: they carry a small power and compute budget and demand results in
real time, but they stress depth differently. A mobile robot needs
reliable obstacle depth at a few meters and tolerates modest resolution;
a drone needs the same at higher frame rates because its closing speeds
are larger; an augmented-reality headset needs sharp depth boundaries so
that virtual content occludes correctly against real edges; and an
embedded sensing rig needs metric accuracy so that the recovered
geometry can be measured and reconstructed. A single network that serves
all four must therefore balance accuracy, boundary sharpness, latency,
and memory rather than excel at any one of them, and that balance is the
design objective of this thesis.

% Eq 1.1: the depth-from-disparity relation (introduce with a purpose
% clause, then a where-clause naming every symbol).
\begin{equation}
Z = \frac{f\,B}{d}
\label{eq:depth_from_disparity}
\end{equation}
where $Z$ is the depth of a scene point, $f$ the camera focal length in
pixels, $B$ the baseline separating the two cameras, and $d$ the
disparity, the horizontal shift of the point between the two rectified
views.
\nm{$d$}{disparity, the horizontal pixel shift between rectified views}
\nm{$Z$}{depth of a scene point along the camera axis}
\nm{$f$}{camera focal length in pixels}
\nm{$B$}{stereo baseline, the distance between the two camera centers}
\nm{$(c_x, c_y)$}{principal point of the camera}
\nm{EPE}{end-point error, the mean absolute disparity error}
\nm{bad-$t$}{percentage of valid pixels with error above $t$ pixels}
\nm{D1-all}{percentage of pixels with error above 3 px and 5\% of ground truth}
\nm{GRU}{gated recurrent unit}
\nm{GEV}{geometry encoding volume}
\nm{KD}{knowledge distillation}
\nm{SGM}{semi-global matching}
\nm{FPS}{frames per second}
\nm{TOPS}{tera-operations per second}
\nm{INT8}{eight-bit integer arithmetic}
\nm{GT}{ground truth}

\begin{figure}[!htbp]
\centering
\includegraphics[width=0.92\textwidth]{fig_1_1_stereo_geometry_preview}
\caption{\textbf{Two-view stereo geometry.} A scene point $P$ projects to $x_L$ and $x_R$ on two rectified image planes separated by the baseline $B$; the disparity $d = x_L - x_R$ recovers depth through $Z = fB/d$.}
\label{fig:stereo_geometry}
\end{figure}

\section{Problem Statement}
% Sources: THESIS_PLAN state-of-play; papers/verified_performance.md
% (FoundationStereo size); review_paper/sections/03_foundation_baseline.tex.
% Content: accurate deep stereo exists (foundation models, hundreds of M
% params) but is undeployable on edge; edge envelope 1 to 6 TOPS, 4 GB,
% 5 to 25 W; accuracy-latency-memory triangle; the concessive gap pivot
% ("However, ...").
Deep networks have transformed the accuracy of stereo matching. The most
accurate systems now pair a foundation-scale monocular prior with an
iterative refinement core; FoundationStereo, for example, reports 0.34
pixel end-point error on the SceneFlow benchmark with a backbone of
roughly 340 million parameters \cite{wen2025foundationstereo}. However,
this accuracy is bought with a computational cost that embedded hardware
cannot pay. The edge accelerators found on mobile robots and drones
deliver a few tera-operations per second, carry a few gigabytes of memory,
and operate within a power envelope of a few watts to a few tens of watts.
A network with hundreds of millions of parameters and hundreds of
milliseconds of latency does not fit that budget. The problem this thesis
addresses is therefore not accuracy alone but the joint objective of
accuracy, latency, and memory: a stereo network that is accurate enough to
be useful, small enough to fit in on-device memory, and fast enough to run
in real time on an edge accelerator. Figure \ref{fig:edge_gap} places the
published literature on this trade-off and marks the region this thesis
targets.

\begin{figure}[!htbp]
\centering
\includegraphics[width=\textwidth]{fig_1_3_edge_gap}
\caption{\textbf{The edge-deployment gap.} SceneFlow end-point error against trainable parameters for representative methods of each family; the shaded band marks the sub-3 M parameter edge budget targeted in this thesis. Published values follow each paper's standard SceneFlow protocol; the StereoLite point is our full FlyingThings3D test result at the native evaluation axis (Chapter 4).}
\label{fig:edge_gap}
\end{figure}

\section{Research Objectives and Scope}
% The two objectives are FIXED by the user (2026-07-05), verbatim below.
% Do NOT reword them. Sub-goals may be listed as scope items.
The objectives of this thesis are:
\begin{enumerate}
\item To design a computationally efficient stereo matching pipeline that
leverages AI-based disparity refinement to enhance depth estimation on
resource-limited platforms.
\item To design an architecture that can withstand camera imperfections
in terms of rectification.
\end{enumerate}

The scope of this thesis is passive rectified binocular stereo, trained by
supervision on synthetic data with dense ground truth. The design targets
a compact network of a few million parameters, evaluated both in domain
and zero-shot on an unseen real domain, with its inference optimized for
the edge latency envelope and qualified on real camera pairs. The
recovered disparity is carried through to metric three-dimensional
reconstruction: real scenes captured by the project's stereo rig are
back-projected into colored point clouds through the calibrated
geometry, which is the reconstruction the thesis title refers to. Active
sensing, self-supervised training, monocular depth estimation, and the
production of a deployment software development kit are outside the scope.
The evidence for the second objective, tolerance to imperfect
rectification, is a controlled perturbation study reported in Chapter 4.

\section{Rationale and Significance}
% Content: why a NEW small architecture rather than shrinking a big one;
% Complex Engineering Problem justification woven in: K3 (engineering
% fundamentals), K4 (specialist stereo knowledge), K5 (network design),
% K8 (research-literature grounding), one sentence each.
A natural question is why this thesis designs a new compact network rather
than compressing an existing large one. Two reasons motivate the choice.
First, the accuracy of the foundation-scale systems rests substantially on
their monocular priors and their scale, neither of which survives
aggressive compression, so compressing them tends to discard the very
component that makes them accurate. Second, the tile-hypothesis and
recurrent-refinement families were designed for efficiency from the
outset and offer a cleaner starting point for an edge budget. Designing
such a network is a complex engineering problem: it cannot be solved by
natural science or mathematics alone. It requires in-depth engineering
knowledge of the stereo pipeline (attribute K3), specialist knowledge of
the deep-stereo literature and its building blocks (K4), the design of a
novel architecture under hard resource constraints (K5), and grounding in
a large body of research literature to select and combine the right
mechanisms (K8). These attributes are demonstrated across the methodology
and results chapters.

\section{Thesis Organization}
% Write LAST for this chapter: one tight sentence per chapter.
The remainder of this thesis is organized as follows. Chapter 2 reviews
the stereo-matching literature, from the classical pipeline through the
deep-learning eras to the modern efficient and foundation-model families,
with an emphasis on how edge deployability is achieved. Chapter 3 presents
the proposed architecture block by block, together with its training
protocol and the optimization pass that meets the edge latency target.
Chapter 4 reports the in-domain and zero-shot accuracy, the ablation
studies, and the efficiency and deployment measurements, and discusses
their implications. Chapter 5 summarizes the findings, states the
limitations, and outlines future work. The appendices provide the
technical specifications, the project management and financial record, the
ethics and teamwork statements, the program-outcome attainment tracker,
and the list of publications.
